% we include the glossary here (frontmatter is included with \input, so this command is as if it was in main.tex)
\input{frontmatter/glossary}

\frontmatter % Use roman page numbering style (i, ii, iii, iv...) for the pre-content pages

\pagestyle{plain} % Default to the plain heading style until the thesis style is called for the body content

%----------------------------------------------------------------------------------------
%	TITLE PAGE
%----------------------------------------------------------------------------------------

\maketitlepage


%----------------------------------------------------------------------------------------
%	STATEMENT of INTEGRITY
%----------------------------------------------------------------------------------------
\integritystatement

%----------------------------------------------------------------------------------------
%	DEDICATION  (optional)
%----------------------------------------------------------------------------------------
%
%\dedicatory{For/Dedicated to/To my\ldots}
\begin{dedicatory}
A dedicatória é opcional.
\end{dedicatory}

%----------------------------------------------------------------------------------------
%	ABSTRACT PAGE
%----------------------------------------------------------------------------------------

\begin{abstract}

% here you put the abstract in the main language of the work.

O resumo do relatório (que só deve ser escrito após o texto principal do relatório estar completo) é uma apresentação abreviada e precisa do trabalho, sem acrescento de interpretação ou crítica, escrita de forma impessoal, podendo ter, por exemplo, as seguintes três partes:

\begin{enumerate}
    \item Um parágrafo inicial de introdução do contexto e do problema/objetivo do trabalho.
    \item Resumo dos aspetos mais importantes do trabalho descrito no presente relatório, que por sua vez documenta abordagem adotada e sistematiza os aspetos relevantes do trabalho realizado durante o estágio. Deve mencionar tudo o que foi feito, por isso deve concentrar-se no que é realmente importante e ajudar o leitor a decidir se quer ou não consultar o restante do relatório.
    \item Um parágrafo final com as conclusões do trabalho realizado.
\end{enumerate}

Palavras-chave (Tema):	Incluir 3 a 6 palavras/expressões chave que caraterizem o projeto do ponto de vista de tema/área de intervenção.

Palavras-chave (Tecnologias): 	Incluir 3 a 6 palavras/expressões chave que caraterizem o projeto do ponto de vista de tecnologias utilizadas.

\vspace{20mm} %5mm vertical space

\textbf{(O Resumo só deve ocupar 1 página, cerca de 20 linhas.)}

\end{abstract}

\begin{abstractotherlanguage}

Here you put the abstract in the "other language": English, if the work is written in Portuguese; Portuguese, if the work is written in English.

\end{abstractotherlanguage}

%----------------------------------------------------------------------------------------
%	ACKNOWLEDGEMENTS (optional)
%----------------------------------------------------------------------------------------

\begin{acknowledgements}

Nesta secção (opcional) colocam-se notas de agradecimento às pessoas que contribuíram para a realização da tarefa. No caso de terem usufruído de financiamento via bolsa ou qualquer outro mecanismo formal, deverá ser incluído um agradecimento.

\end{acknowledgements}


%----------------------------------------------------------------------------------------
%	LIST OF CONTENTS/FIGURES/TABLES PAGES
%----------------------------------------------------------------------------------------

\tableofcontents % Prints the main table of contents

\listoffigures % Prints the list of figures

\listoftables % Prints the list of tables

\iflanguage{portuguese}{
	\renewcommand{\listalgorithmname}{Lista de Algor\'itmos}
}
\listofalgorithms % Prints the list of algorithms
%\addchaptertocentry{\listalgorithmname}


\renewcommand{\lstlistlistingname}{List of Source Code}
\iflanguage{portuguese}{
	\renewcommand{\lstlistlistingname}{Lista de C\'odigo}
}
\lstlistoflistings % Prints the list of listings (programming language source code)
%\addchaptertocentry{\lstlistlistingname}


%----------------------------------------------------------------------------------------
%	ABBREVIATIONS
%----------------------------------------------------------------------------------------
%\begin{abbreviations}{ll} % Include a list of abbreviations (a table of two columns)
%%\textbf{LAH} & \textbf{L}ist \textbf{A}bbreviations \textbf{H}ere\\
%%\textbf{WSF} & \textbf{W}hat (it) \textbf{S}tands \textbf{F}or\\
%\end{abbreviations}

%----------------------------------------------------------------------------------------
%	SYMBOLS
%----------------------------------------------------------------------------------------

%\begin{symbols}{lll} % Include a list of Symbols (a three column table)

%$a$ & distance & \si{\meter} \\
%$P$ & power & \si{\watt} (\si{\joule\per\second}) \\
%%Symbol & Name & Unit \\

%\addlinespace % Gap to separate the Roman symbols from the Greek

%$\omega$ & angular frequency & \si{\radian} \\

%\end{symbols}

%----------------------------------------------------------------------------------------
%	ACRONYMS
%----------------------------------------------------------------------------------------

\newcommand{\listacronymname}{List of Acronyms}
\iflanguage{portuguese}{
	\renewcommand{\listacronymname}{Lista de Acr\'onimos}
}

%Use GLS
\glsresetall
\printglossary[title=\listacronymname,type=\acronymtype,style=long]

%----------------------------------------------------------------------------------------
%	DONE
%----------------------------------------------------------------------------------------

\mainmatter % Begin numeric (1,2,3...) page numbering
\pagestyle{thesis} % Return the page headers back to the "thesis" style
